% Draft status: V0
% Evidence status: checked where possible; TODOs mark missing evidence.
\section{Analysis and Discussion}

The results suggest a trade-off-like pattern, but the draft should describe it cautiously. Finetuning is associated with much higher object accuracy for every model-dataset pair, and it is also associated with a larger regional bias gap for every model-dataset pair. One possible explanation is that the pretrained evaluation has uniformly weak performance, which compresses regional differences because most predictions are incorrect. After finetuning, models become more capable of recognizing the target classes, and differences between regions become more visible. This explanation is plausible, but it cannot establish causality and requires more evidence about training dynamics, class distributions, and per-region sample sizes.

The Dollar Street results are especially important for the main research question. Every finetuned Dollar Street model has Europe as the best region and Africa as the worst region. This is consistent with prior work suggesting that object recognition can perform unevenly across countries and income contexts in Dollar Street-like imagery \citep{devries2019object}. However, the current paper should not claim the same mechanism or the same demographic interpretation without further evidence. The observed pattern may reflect dataset composition, object categories, image conditions, socioeconomic context, model pretraining data, or interactions among these factors.

GeoDE shows a different pattern. Finetuned object accuracies are high across all nine models, and regional gaps are smaller than on Dollar Street. Most finetuned GeoDE models have WestAsia as the best region and SouthEastAsia as the worst region. This suggests that GeoDE still reveals regional variation even when overall performance is high. At the same time, the lower gap values indicate that the regional disparities measured by this paper's metric are smaller on GeoDE than on Dollar Street. TODO: evidence needed to explain whether GeoDE's dataset construction or class distribution is associated with this smaller measured gap.

High object accuracy is not always associated with low regional bias gap. ConvNeXt has the highest finetuned object accuracy on both datasets: 0.7667 on Dollar Street and 0.9661 on GeoDE. Its finetuned regional gap is 0.0848 on Dollar Street and 0.0229 on GeoDE. MAE has lower finetuned object accuracy than ConvNeXt on both datasets, but it has the lowest finetuned gap on both datasets: 0.0688 on Dollar Street and 0.0187 on GeoDE. This suggests that accuracy and regional stability are not the same empirical property. This is consistent with worst-group evaluation arguments that aggregate accuracy can obscure group-specific performance differences \citep{sagawa2020groupdro}.

The model-family comparison is suggestive but incomplete. In a rough grouping, ConvNeXt and ResNet represent CNN models; DeiT, Swin, and MaxViT represent transformer models; DINO and MAE represent self-supervised models; and CLIP and SigLIP represent vision-language models. Averaging from the normalized table, the finetuned transformer group has mean object accuracy 0.8537 and mean regional gap 0.0554, while the finetuned vision-language group has mean object accuracy 0.8514 and mean regional gap 0.0584. The self-supervised group has lower mean finetuned object accuracy at 0.8224 but also the lowest mean finetuned regional gap at 0.0471. The CNN group has mean finetuned object accuracy 0.8075 and mean finetuned regional gap 0.0641. These family-level numbers are useful as a first draft summary, but a dedicated family-level table should be added before making strong claims.

The vision-language results are mixed. CLIP has high finetuned GeoDE accuracy but the largest finetuned GeoDE gap, while SigLIP is among the stronger Dollar Street finetuned models by object accuracy and has a lower Dollar Street gap than CLIP. Prior CLIP robustness work suggests that properties of pretraining data may be associated with distributional robustness \citep{fang2022cliprobustness}, but this paper cannot test that mechanism directly. The safest interpretation is that vision-language models in this benchmark do not uniformly dominate on both object accuracy and regional gap.

The self-supervised results are also mixed. MAE has the lowest finetuned regional gap on both datasets but lower object accuracy than several alternatives, especially on Dollar Street. DINO has higher Dollar Street finetuned object accuracy than MAE but a larger Dollar Street regional gap. DINOv2 and MAE are both commonly framed as self-supervised representation-learning approaches \citep{oquab2024dinov2,he2022mae}, but the current results do not support a broad claim that self-supervision automatically improves regional robustness.

The gap increase after finetuning should be interpreted carefully. It may indicate that better class discrimination exposes regional differences; it may reflect uneven learning across regions; or it may be partly a mathematical consequence of moving from near-zero accuracy to high accuracy. These possibilities are not mutually exclusive. Additional analyses, such as per-class regional gaps, region sample sizes, confidence calibration, and training curves, would help distinguish them. TODO: evidence needed for per-class regional disparity and uncertainty estimates.

Overall, the results support the paper's main premise: aggregate object accuracy is not enough to characterize regional model behavior. The same model can look strong by object accuracy and still show a sizable best-worst regional gap. Conversely, a low regional gap in a very low-accuracy setting should not be interpreted as robust performance. This reinforces the need to report object accuracy, worst-region accuracy, best-region accuracy, and regional bias gap together.

